\section{Discusión}
El análisis comparativo entre los resultados del proyecto y los hallazgos de la literatura 
evidencia coincidencias y aprendizajes relevantes. En primer lugar, la selección de 
Laravel como framework backend se justifica no solo por su facilidad de uso, sino 
también porque los estudios demuestran su capacidad para garantizar sistemas robustos y 
escalables en entornos de producción [1]. En comparación con Symfony, Laravel 
presenta una curva de aprendizaje más corta, lo que lo convierte en la opción ideal para 
equipos pequeños o en formación [7]. 

En segundo lugar, el uso de React permitió construir una interfaz más moderna y basada 
en componentes reutilizables, lo que coincide con estudios que resaltan su liderazgo en 
proyectos de gran escala [2], [4], [5]. La incorporación de React Native se vislumbra 
como una evolución natural del proyecto, en línea con experiencias previas que destacan 
su eficacia en el desarrollo multiplataforma [3]. 

Respecto a las metodologías ágiles, el marco de trabajo Scrum resultó útil para mantener 
el ritmo del proyecto; no obstante, como lo señalan autores expertos, su uso debe ir 
acompañado de prácticas robustas de documentación, arquitectura y pruebas constantes 
[6], [9]. Esta combinación fue fundamental para disminuir la deuda técnica y garantizar 
que el software tenga calidad. 

Finalmente, en el proyecto se confirma que la integración de una teoría y práctica, se  
fortalece en el desarrollo de soluciones tecnológicas. Los artículos revisados no solo 
ayudaron como referencias bibliográficas, sino que guiaron las decisiones estratégicas 
durante el proceso de implementación, lo cual se tradujo en un producto funcional, 
seguro y alineado con las tendencias actuales del desarrollo de software. 